% !PRINTED 485
\weberpage{485}
%<<CONT
und wenn man die beiden Gleichungen (35) addiert:
\[
\begin{aligned}
&[(\xi+\eta)^3-3\,\xi\eta(\xi+\eta)]\left(1+\frac{1}{\xi^3\eta^3}\right)\\
&+\,2[(\xi+\eta)^2-2\,\xi\eta]\left(1-\frac{1}{\xi^2\eta^2}\right)=0.
\end{aligned}
\]

Es ist aber nach (34):
\[
\sqrt2\,\xi\eta=x^2,
\]
und man erhält also eine Gleichung für $x^2$, wenn man $(\xi+\eta)$
eliminiert. Diese Gleichung wird nach dem gewöhnlichen Eliminationsverfahren, wenn man
\[
z=\frac{x^2}{\sqrt2}+\frac{\sqrt2}{x^2}
\]
setzt, in der Form
\weq{36}{z^6-9\,z^5+20\,z^4+6\,z^3-19\,z^2-17\,z-6=0}
gefunden. Hierin ist aber noch der auf die Determinante $-49$ bezügliche
Faktor enthalten. Für diesen ist [nach Formel (10)]:
\[
z=2+\sqrt7.
\]

Es muß also die linke Seite von (36) durch
\[
z^2-4\,z-3
\]
teilbar sein, und die Ausführung der Division ergibt den für die
Determinante $-41$ gültigen Faktor:
\weq{37}{\begin{aligned}
z^4-5\,z^3+3\,z^2+3\,z+2&=0, \\
z&=\dfrac{f(\sqrt{-41})}{\sqrt2}+\dfrac{\sqrt2}{f(\sqrt{-41})}.
\end{aligned}}

\begin{center}
\S\,131. \gesperrt{Berechnung von Klasseninvarianten aus den irrationalen Formen der Modulargleichungen.}
\end{center}

In außerordentlich einfacher Weise führen vielfach die irrationalen
Formen der Modulargleichungen zur Aufstellung von Klassengleichungen.

1. Im \S\,75 haben wir gesehen, daß, falls $m\equiv -1\ (\mathrm{mod}\ 8)$
ist, zwischen den beiden Funktionen:
\[
\begin{aligned}
2A&=f(\omega)f(m\omega)+(-1)^{\frac{m+1}{8}}[f_1(\omega)f_1(m\omega)+f_2(\omega)f_2(m\omega)],\\
B&=\frac{2}{f_1(\omega)f_1(m\omega)}+\frac{2}{f_2(\omega)f_2(m\omega)}+\frac{(-1)^{\frac{m+1}{8}}2}{f(\omega)f(m\omega)}
\end{aligned}
\]
%CONT>>
